\chapter{内存层级与算子实现}\label{v0ch:26}
\begin{chaptersummary}
算子的理想流量给出了复用目标，实际实现还需确定数据经过哪些存储层。本章将寄存器、缓存和 DRAM 联系到布局、分块、融合与 SIMD，逐项分析临时量、依赖和访问路径。有效优化由此落实为减少关键层级的数据移动，并在同一误差与计时口径下检验收益。
\end{chaptersummary}
\learninggoals{分析步长与缓存复用；解释融合、分块、SIMD 和矩阵自由的取舍；设计公平计时。}
\paragraph{核心问题}理论上只需读一次的数据，程序为什么会反复搬运？

\section{层级、局部性与布局}
寄存器靠近执行单元，L1/L2/L3 缓存逐级扩大容量，DRAM 提供主要容量。具体共享方式依架构而异。连续访问有利于缓存行利用和硬件预取，跨大步长或随机 gather 可能浪费事务；工作集尺寸决定哪些复用能真正保留。
C 顺序三维数组最后一个索引连续；循环与线程映射宜使连续维成为内层。多分量 SoA 常利于跨点向量化，AoS 可能利于单点全部分量复用，AoSoA 则在小块内折中。应以实际核的数据使用方式判断。
\section{融合与分块}
表达式 $w=a u+b v$ 若先形成两个临时数组再相加，会增加写回和读取；融合可直接读 $u,v$ 写 $w$。但融合增加寄存器活跃变量，过度融合可能溢出至内存。
空间分块将 stencil 相邻平面留在缓存；时间分块跨多个时间步复用子域，必须保留依赖锥和扩大 halo。Yee 的两相更新或多阶段 RK 不能在缺少同步的条件下任意跨步融合。
张量 DG 可将导数矩阵留在缓存并批量处理单元；若同时保存三方向导数和通量，临时量可能超过寄存器或缓存。矩阵自由减少稀疏索引但增加重复几何计算，混合精度还需误差论证。
\conceptbridge{“零拷贝”或“矩阵自由”为何仍可能很慢？}{避免显式复制或矩阵条目只是减少某一类操作。缓存失配、隐式临时量、跨层迁移和额外算术仍可主导时间。应依据实际访问路径，而非接口名称判断数据流量。}{存储表示 $\longrightarrow$ 实际访问路径 $\longrightarrow$ 各层资源消耗。}
\section{SIMD 与最小实验}
SIMD 一条指令处理多个数据元素，要求独立迭代、合适布局、对齐及可处理的分支。编译器向量化报告、汇编或计数器比源代码外观更可靠。
对七点 stencil 比较连续内层与跨步长遍历，对 D2Q9 比较 SoA/AoS，对张量导数比较显式 Kronecker 与顺序收缩。计时前预热，保留输出以防消除，重复并报告统计值；大数组和缓存驻留小数组分开。
\costnote{分别报告分配时间与核函数时间，在统一计时范围下比较预分配和动态分配实现。NumPy 的各数组操作调用编译核，多操作之间的融合则取决于具体执行方式。}
\coderef{roofline/kernels.py; lbm/stream_collision.py; dg_fr/transport.py}
\aiexercise{逐数组列出 NumPy 七点 stencil 的临时量与生命周期，比较其字节数和第 23 章融合核下界。}{设计三个工作集分别小于 L1、L3 和大于末级缓存的实验，解释为什么同一强度点不能共用一个未经区分的带宽。}
\studyroute{略去缓存替换策略、编译器向量化与具体微架构端口分析。参考 \cite{roofline,cuda}。}{cache blocking; loop fusion; temporal blocking; SIMD; AoSoA; working set}
